\section[Lý thuyết]{Mở đầu và lý thuyết rút gọn}

\begin{frame}{Xây kho dữ liệu và BI cho 99.441 đơn hàng thương mại điện tử Olist}
\begin{columns}[T]\begin{column}{0.4\textwidth}
\begin{itemize}
  \item Học phần MI4214 – Kho dữ liệu và Kinh doanh thông minh; nhóm ...; giảng viên ...
\end{itemize}
\end{column}\begin{column}{0.58\textwidth}
\begin{center}\anh[0.74\textheight]{3-3-4_constellation.png}\end{center}
\end{column}\end{columns}

\note{[S01] Chúng em xây một kho dữ liệu hoàn chỉnh từ hai bộ dữ liệu Olist: nguồn, ETL, kho hình sao, data mart và bảy dashboard Power BI. Bài trình bày đi từ bài toán đến kết quả và bài học.}
\end{frame}

\begin{frame}{Bài trình bày gồm bảy phần, đi từ lý thuyết đến sản phẩm chạy được}
\begin{itemize}
  \item 1 Lý thuyết $\cdot$ 2 Bài toán và nghiệp vụ $\cdot$ 3 Dữ liệu và khám phá $\cdot$ 4 Kiến trúc và thiết kế $\cdot$ 5 ETL, mart, OLAP, hiệu năng $\cdot$ 6 Dashboard và mở rộng $\cdot$ 7 Kết luận
\end{itemize}


\note{[S02] Mỗi phần có số liệu thật từ kho; cuối mỗi phần nhóm nêu điều học được.}
\end{frame}

\begin{frame}{Dự án trong một trang: 1,56 triệu dòng nguồn thành 6 fact, 17 view mart và 7 dashboard}
\begin{columns}[T]\begin{column}{0.4\textwidth}
\begin{itemize}
  \item 2 hệ thống nguồn $\rightarrow$ staging $\rightarrow$ ODS $\rightarrow$ kho (10 chiều + 6 fact) $\rightarrow$ 4 data mart $\rightarrow$ Power BI
\end{itemize}
\end{column}\begin{column}{0.58\textwidth}
\begin{center}\anh[0.74\textheight]{3-4-1_pipeline.png}\end{center}
\end{column}\end{columns}

\note{[S03] Toàn bộ dựng lại từ đầu bằng một lệnh trong khoảng 7 phút và cho kết quả giống hệt; 136 phép kiểm thử đạt. Đây là bức tranh để quay lại khi cần.}
\end{frame}

\begin{frame}{Kho dữ liệu là nơi tích hợp dữ liệu hướng chủ đề, biến thiên theo thời gian và không bị ghi đè}
\begin{itemize}
  \item Bốn đặc tính của Inmon: hướng chủ đề, tích hợp, biến thiên theo thời gian, không biến động
\end{itemize}
\khungchen[0.3\textheight]{CHÈN HÌNH / BẢNG: sơ đồ 4 ô (vẽ)}

\note{[S04] Ở dự án này chủ đề là quản lý đơn hàng; tích hợp hai nguồn bán hàng và marketing; lịch sử khách được giữ bằng SCD2.}
\end{frame}

\begin{frame}{OLTP tối ưu ghi từng giao dịch, OLAP tối ưu đọc và tổng hợp}
\begin{itemize}
  \item OLTP: 9 bảng chuẩn hóa; OLAP: một fact nối các chiều
\end{itemize}
\begin{columns}[c]\begin{column}{0.3\textwidth}\anh[0.5\textheight]{3-2-4_erd_oltp.png}\end{column}\begin{column}{0.68\textwidth}\begin{center}\adjustbox{max width=\linewidth,max height=0.5\textheight}{\scriptsize\renewcommand{\arraystretch}{1.15}\begin{tabular}{lll}\toprule
\textbf{Tiêu chí} & \textbf{OLTP} & \textbf{OLAP (kho dữ liệu)} \\\midrule
Mục đích & ghi từng giao dịch & phân tích, tổng hợp \\
Mô hình dữ liệu & 9 bảng chuẩn hóa (\texttt{oltp}) & hình sao: 6 fact + 10 chiều (\texttt{dw}) \\
Truy vấn điển hình & nối nhiều bảng, ít dòng & quét và gom nhiều dòng \\
Dữ liệu & trạng thái hiện tại & có lịch sử (SCD2 cho khách) \\
Ví dụ & \texttt{orders}, \texttt{order\_items} & \texttt{fact\_order\_item}, \texttt{dim\_date} \\
\bottomrule\end{tabular}}\end{center}\end{column}\end{columns}

\note{[S05] Cùng một câu hỏi "GMV theo bang theo tháng" cần nối ba bảng ở nguồn nhưng chỉ nối một fact với hai chiều ở kho.}
\end{frame}

\begin{frame}{Kimball xây từ các data mart có chiều dùng chung, Inmon xây kho 3NF rồi mới tách mart}
\begin{itemize}
  \item Nhóm chọn hướng Kimball: bus matrix, conformed dimension
\end{itemize}
\begin{center}\adjustbox{max width=\linewidth,max height=0.5\textheight}{\scriptsize\renewcommand{\arraystretch}{1.15}\begin{tabular}{lll}\toprule
\textbf{Tiêu chí} & \textbf{Kimball} & \textbf{Inmon} \\\midrule
Cách xây & từ data mart theo quy trình nghiệp vụ, chiều dùng chung & kho doanh nghiệp chuẩn hóa 3NF, rồi tách data mart \\
Mô hình & chiều (lược đồ sao) & chuẩn hóa \\
Thời gian ra kết quả & nhanh, từng quy trình một & lâu hơn, thiết kế tổng thể trước \\
Nhóm chọn & \ding{51} (bus matrix, \texttt{dim\_seller} dùng chung) & -- \\
\bottomrule\end{tabular}}\end{center}

\note{[S06] Vì bài toán có hai quy trình nghiệp vụ dùng chung chiều seller và ngày, cách Kimball phù hợp và nhanh ra kết quả.}
\end{frame}

\begin{frame}{Mô hình chiều gồm fact, dimension và grain: chọn grain trước, mọi thứ còn lại theo sau}
\begin{columns}[T]\begin{column}{0.4\textwidth}
\begin{itemize}
  \item Grain của dự án: 1 dòng hàng, 1 đơn, 1 lần thanh toán, 1 review, 1 MQL, 1 seller $\times$ tháng
\end{itemize}
\end{column}\begin{column}{0.58\textwidth}
\begin{center}\anh[0.74\textheight]{3-3-4_erd_dw.png}\end{center}
\end{column}\end{columns}

\note{[S07] Grain mịn nhất cho phép cộng lên mọi cấp; ví dụ GMV tính từ grain dòng hàng.}
\end{frame}

\begin{frame}{Ba loại fact và hai kiểu SCD đủ để mô tả mọi nhu cầu của bài toán}
\begin{itemize}
  \item Transaction $\cdot$ periodic snapshot $\cdot$ accumulating snapshot; SCD1 (ghi đè) và SCD2 (giữ lịch sử)
\end{itemize}
\begin{center}\adjustbox{max width=\linewidth,max height=0.5\textheight}{\scriptsize\renewcommand{\arraystretch}{1.15}\begin{tabular}{lll}\toprule
\textbf{Loại} & \textbf{Một dòng là} & \textbf{Ví dụ trong dự án} \\\midrule
Transaction & một sự kiện đã xảy ra & \texttt{fact\_order\_item}, \texttt{fact\_payment}, \texttt{fact\_review} \\
Periodic snapshot & trạng thái của một kỳ & \texttt{fact\_seller\_monthly} (36.583 dòng, dày đặc) \\
Accumulating snapshot & cả vòng đời, cập nhật theo mốc & \texttt{fact\_order\_fulfillment} (99.441), \texttt{fact\_lead\_funnel} (8.000) \\
SCD1 / SCD2 & ghi đè / giữ lịch sử & \texttt{dim\_seller} / \texttt{dim\_customer} \\
\bottomrule\end{tabular}}\end{center}

\note{[S08] Dự án có đủ cả ba loại fact và dùng SCD2 cho khách, SCD1 cho seller.}
\end{frame}

\begin{frame}{ETL biến dữ liệu nguồn bẩn thành dữ liệu đáng tin qua staging và ODS}
\begin{columns}[T]\begin{column}{0.4\textwidth}
\begin{itemize}
  \item Extract $\rightarrow$ làm sạch ở ODS $\rightarrow$ nạp chiều và fact; full và incremental
\end{itemize}
\end{column}\begin{column}{0.58\textwidth}
\begin{center}\anh[0.74\textheight]{3-4-1_pipeline.png}\end{center}
\end{column}\end{columns}

\note{[S09] Nguyên tắc của nhóm là không xóa âm thầm: dòng lỗi được gắn cờ hoặc ghi vào bảng loại kèm lý do.}
\end{frame}

\begin{frame}{Năm thao tác OLAP đều quy về SQL trên lược đồ sao}
\begin{itemize}
  \item Roll-up, drill-down, slice, dice, pivot
\end{itemize}
\begin{center}\adjustbox{max width=\linewidth,max height=0.5\textheight}{\scriptsize\renewcommand{\arraystretch}{1.15}\begin{tabular}{lll}\toprule
\textbf{Thao tác} & \textbf{SQL} & \textbf{Ví dụ trong dự án} \\\midrule
Roll-up & \texttt{ROLLUP} & quý $\rightarrow$ năm $\rightarrow$ tổng \\
Drill-down & thêm cấp vào \texttt{GROUP BY} & 2017 $\rightarrow$ Q4 $\rightarrow$ 11/2017 $\rightarrow$ ngày \\
Slice & \texttt{WHERE} một chiều & năm 2017 \\
Dice & \texttt{WHERE} nhiều chiều & Sudeste $\times$ Home \& Furniture $\times$ Q4 \\
Pivot & \texttt{FILTER} & vùng $\times$ năm \\
\bottomrule\end{tabular}}\end{center}

\note{[S10] Mỗi thao tác là một câu SQL ngắn; phần 5 sẽ cho kết quả thật.}
\end{frame}

\begin{frame}{MOLAP tính sẵn khối đa chiều, ROLAP tính lúc truy vấn, HOLAP kết hợp: nhóm chọn ROLAP}
\begin{itemize}
  \item Lý do: PostgreSQL không có khối đa chiều, dữ liệu 112.650 dòng hàng trả lời dưới 1 giây, kho luôn cập nhật sau mỗi lần nạp
\end{itemize}
\begin{center}\adjustbox{max width=\linewidth,max height=0.5\textheight}{\scriptsize\renewcommand{\arraystretch}{1.15}\begin{tabular}{llll}\toprule
\textbf{Tiêu chí} & \textbf{MOLAP} & \textbf{ROLAP} & \textbf{HOLAP} \\\midrule
Lưu trữ & khối đa chiều riêng & bảng quan hệ (lược đồ sao) & kết hợp hai cách \\
Tốc độ truy vấn & rất nhanh (tính sẵn) & phụ thuộc chỉ mục, materialized view & cân bằng \\
Độ tươi dữ liệu & phải xử lý lại khối & luôn tươi sau mỗi lần nạp & tùy phần \\
Công cụ & ví dụ SSAS & PostgreSQL + Power BI & -- \\
Nhóm chọn & -- & \ding{51} ROLAP & -- \\
\bottomrule\end{tabular}}\end{center}

\note{[S11] Đổi lại, truy vấn nặng được tăng tốc bằng materialized view; phần hiệu năng sẽ cho thấy lợi ích.}
\end{frame}

\begin{frame}{BI đưa kho dữ liệu đến người ra quyết định qua mô hình ngữ nghĩa, KPI và dashboard}
\begin{itemize}
  \item Power BI: nạp CSV, quan hệ, measure DAX, dashboard
\end{itemize}
\khungchen[0.3\textheight]{CHÈN HÌNH / BẢNG: sơ đồ chuỗi dữ liệu $\rightarrow$ kho $\rightarrow$ mô hình $\rightarrow$ dashboard (vẽ)}

\note{[S12] Mọi KPI được định nghĩa một lần bằng DAX và đối chiếu với giá trị tính bằng SQL.}
\end{frame}
